\documentclass[11pt,oneside]{book}
\pagestyle{plain}
\usepackage{graphics}
\usepackage{makeidx}
\usepackage{nomencl}

\textheight=235mm
\textwidth=145mm
\topmargin=0mm
\headheight=0mm
\headsep=0mm

% These must be changed for double sided output !

\oddsidemargin=15mm
\evensidemargin=15mm

\makeindex
\makeglossary
\begin{document}

% set line spacing to 1.5B
\baselineskip = 17pt


\title{CSD \LaTeX \ Estimación de la mirada}
\author{Thesis submitted in accordance with the requirements of \\
the University of Liverpool for the degree of Doctor in Philosophy \\
by \\
Ian C. Smith}
\date{December 2005}
\maketitle 
\frontmatter

\chapter{Resumen}


\tableofcontents

% Contents not put in table of contents by default so add it separately

\addcontentsline{toc}{chapter}{Contents}

\listoffigures


% List of figures  not put in table of contents by default so add it separately

\addcontentsline{toc}{chapter}{List of Figures}


\chapter{Acknowledgement}

The author wishes to thank the University of
Liverpool Computing Services Department for the development of this
\LaTeX \ thesis template.

\printglossary

\addcontentsline{toc}{chapter}{Nomenclature}





\mainmatter

\chapter{How to use the Thesis Template}
\section{Introduction}

\LaTeX \ is a document publishing system which can produce high quality 
output particulary when complicated formulae need to be reproduced in print.
Because of this it is widely used by mathematicians. \LaTeX \ is essentially
a collection of \TeX \ macros which is a typesetting language in its own right.
\TeX \ is quite difficult to learn but provides much finer control over typsetting
then \LaTeX. Most users will find that \LaTeX \ meets most of their
needs and will therefore have little need of \TeX. \LaTeX \ is available on 
both the MWS and Sun UNIX services with the MWS
version having the added advantage of including a graphical user interface called 
TeXShell.

There are no courses provided by CSD in \LaTeX \ however, information is available
on line. The TeXShell/MikTeX guide can be found at:
\\
\\
http://www.liv.ac.uk/CSD/acuk\_html/414.dir/414.doc
\\
\\
and there is also an FAQ at:
\\
\\
http://www.liv.ac.uk/csd/software/ps/faqtex.htm
\\
\\
The Sun UNIX \LaTeX \ guide is at:
\\
\\
http://www.liv.ac.uk/CSD/acuk\_html/452.dir/452.pdf
\\
\\
If you are new to \LaTeX, then probably the best place to start is with 
Leslie Lamport's very readable guide, {\em ``LaTeX - A Document Preparation
System''} \cite{latex} . If you need to delve into \TeX, then {\em ``The TeXbook''} \cite{texbook} by 
Donald E. Knuth is essential reading. Knuth, a leading light in 
computer science, designed the \TeX \ language in order to produce
his seminal work {\em ``The Art of Computer Programming''} \cite{artbook} so this really
is the definitive guide to the subject.

To save users from having to type in a many low level
\TeX \ commands, style files can be used to define additional
macros which can be used in addition to standard ones from
\LaTeX. Many journal publishers provide style files to ensure
that all submissions conform to a standard `house' format. It may
be worth checking on the Internet to see if style files are available
to suit your particular discipline.
 
\section{Running \LaTeX}
This section descrines how to take the thesis template file
({\bf thesis.tex}) and convert it into a PostScript or Portable
Document Format (PDF) file which can be printed on a high quality
laser printer. Both PostScript and PDF files can be previewed
prior to printing using either Ghostview or Adobe's Acorobat
viewer. There are some common pitfalls and misconceptions associated with
with Ghostview and PostScript, and neophyte\nomenclature{{\bf Neophyte}}{A beginner or novice.} users are strongly
advised to take a look at the ``FAQ for Ghostview and PostScript"
at:
\\
\\
http://www.liv.ac.uk/csd/software/latex/faqtex.htm
\\
\\
Both \LaTeX \ and \TeX \ convert the original ``source" file
into a DVI file (by convention having the {\bf .dvi} extension)
before this is converted to PostScript. DVI is short
for DeVice Independent format, meaning that it works
independently of the printer format (usually PostScript)
which is used to produced the hard copy version. Despite
what the name suggests, DVI files are seldom useful outside
of the \LaTeX \ environment.

Historically speaking, PostScript has been widely used to 
publish \LaTeX \ derived articles
on the World Wide Web but, since the arrival of PDF, this should be
stricly discouraged. PostScript is a format designed for use with
laser printers and many modern web browsers have difficulty displaying
files which contain it. On the other hand most browsers will display
PDF files and PDF is becoming an increasing popular standard for sharing
textual information
\footnote{
You might wish to consider how best to save your thesis in electronic format
for posterity. Saving the \LaTeX \ original as a plain text file and 
the formatted version as a PDF file
onto 
ordinary CD-R is probably the safest bet at the moment, although bear in mind 
that the computer graveyard is full of ``standards" of depressingly short longevity.
Fortunately, \LaTeX \ is likely to be far more future proof than this week's version
of Microsoft Word.
}.

\subsection{\LaTeX \ on the Sun UNIX Service}

To ``compile" the template file ({\bf thesis.tex}) into a DIV file, the
following command is used:

\begin{verbatim}
$ latex thesis.tex   
\end{verbatim}

Alternatively, in the unlikely circumstances that the input file contains pure
\TeX \ , use: 

\begin{verbatim}
$ tex pure_tex_file.tex   
\end{verbatim}

Assuming that this does not throw up any errors, you can then convert the
DVI file to PostScript format using:

\begin{verbatim}
$ dvips thesis  
\end{verbatim}

You can then preview the output to ensure that what you get from the
printer is actually what you wanted using Ghostview as 
follows:

\begin{verbatim}
$ ghostview thesis.ps &
\end{verbatim}


Finally, to print the output to a laser printer, use:

\begin{verbatim}
$ lpr -Pps thesis.ps
\end{verbatim}


By default, this will go to the central printer in Brownlow Hill;
the {\bf -P} option can be used to direct the output elsewhere. It is 
straightforward to convert the PostScript version to PDF:

\begin{verbatim}
$ ps2pdf thesis.ps thesis.pdf
\end{verbatim}

and infact it is possible to convert \LaTeX \ documents to many other formats, including
HTML. See the TeX FAQ for details.


\subsection{\LaTeX \ on the Managed Windows Service}
\LaTeX \ on the Managed Windows Service comes in two parts namely:
a group of programs run from the DOS command line called MikTeX
and a graphical windows interface called TeXShell. If neither
of these appear to be installed on your MWS PC click

\begin{verbatim}
Start | Install | Office | MikTeX 2.1
\end{verbatim}

The TeXShell interface is by far the easiest to work with although
MikTeX provides additional programs not accessible through
the standard TeXShell interface
\footnote
{Although they can be included --- see the FAQ for details.}.
To compile the template file click on 
\begin{verbatim}
File | Open
\end{verbatim}
in TeXShell and browse to where you have saved {\bf thesis.tex}.
Then click on the LaTeX button (use the TeX button for plain \TeX \ files).

To convert the DVI output
to PostScript, click on the Dvips button. You can then 
preview the output prior to printing by clicking the
Ghostview button. The document can then be printed
by clicking File $|$ Print inside Ghostview.

\section{Some examples of \LaTeX \ text formatting}

\subsection {Plain Text}
Type your text in free-format; lines can be as long
or as short
as you wish.
        You can indent         or space out
        your input 
            text in 
                any way you like to highlight the structure
        of your manuscript and make it easier to edit.
LaTeX fills lines and adjusts spacing between words to produce an
aesthetically pleasing result.

Completely blank lines in the input file break your text into
paragraphs.
To change the font for a single character, word, or set of words, 
enclose the word and the font changing command within braces, 
{\em like this}.
A font changing command not enclosed in braces, like the change to \bf 
bold here, keeps that change in effect until the end of the document or
until countermanded by another font switch, like this change back to 
\rm roman.

\subsection {Displayed Text}
Use the ``quote'' and  ``quotation'' environments for typesetting quoted
material or any other text that should be slightly indented and set off
from the normal text.
\begin{quotation}
The quote and quotation environments are similar, but use different 
settings for paragraph indentation and spacing.

\em When in doubt, consult the manual.
\end{quotation}


\begin{enumerate}
\item 
The ``enumerate'' environment numbers the list elements, like this.

Items in a list can contain multiple paragraphs.
These paragraphs are appropriately spaced and indented according to their
position in the list.
   \begin{itemize}
   \item The ``itemize'' environment sets off list items with ``bullets'',
like this.  Finally, the ``description'' environment lets you put your own
      \begin{description}
      \item[A] label on each item, like this ``A''.
      \item[If the label is long,] the first line of the item text will
be spaced over to the right as needed.
      \end{description}
   \item Of course, lists can be nested, each type up to at least four levels.
One type of list can be nested within another type.
      \begin{itemize}
      \item Nested lists of the same type will change style of numbering 
or ``bullets'' as needed.
      \end{itemize}
   \end{itemize}
\item Don't forget to close off all list environments with the 
appropriate \verb+\end{...}+ command.
Indenting \verb+\begin{...}+, \verb+\item+, and \verb+\end{...}+
commands in the input document according to their nesting level can help 
clarify the structure.
\end{enumerate}

Here is a very simple table showing data lined up in columns.
Notice that the table is in a ``center'' environment to display
it properly.
The title is created simply as another paragraph in the center environment,
rather than as part of the table itself.
\begin{center}
Numbers of Computers Network, By Type.

\begin{tabular}{lr}
Macintosh&175\\
DOS/Windows PC&60\\
UNIX Workstation or server&110\\
\end{tabular}
\end{center}

Here is a more complicated table that has been boxed up, with a multi-column
header and paragraph entries set in one of the columns.
\begin{center}
\begin{tabular}{|l|c|p{3.5in}|}
\hline
\multicolumn{3}{|c|}{Places to Go Backpacking}\\ \hline
Name&Driving Time&Notes\\
&(hours)&\\ \hline
Big Basin&1.5&Very nice overnight to Berry Creek Falls from
either Headquarters or ocean side.\\ \hline
Sunol&1&Technicolor green in the spring.  Watch out for the cows.\\ \hline
Henry Coe&1.5&Large wilderness nearby suitable for multi-day treks.\\ \hline
\end{tabular}
\end{center}

\subsection {Mathematical Equations}
Simple equations, like $x^y$ or $x_n = \sqrt{a + b}$ can be typeset right
in the text line by enclosing them in a pair of single dollar sign symbols.

A more complicated equation should be typeset in {\em displayed math\/} mode,
like this:
\[
p(x)=\lim_{N \rightarrow \infty, \Delta x \rightarrow 0}\sum_{j=1}^{K}(f_j/\Delta x)
\]
The ``equation'' environment displays your equations, and automatically
numbers them consecutively within your document, like this:
\begin{equation}
\mu_m=\sum_{i=0}^m(-1)^i{m \choose i}\mu_1^i\mu_{m-i}^{\prime} \; \; {\rm where} \; \; {m \choose i}={m! \over i!(m-i)!}.
\label{eq:summation1}
\end{equation}

\pagebreak
\noindent{\Large\bf Here is the input file that produced these formatting examples:}
\begin{verbatim}

\subsection {Plain Text}
Type your text in free-format; lines can be as long
or as short
as you wish.
        You can indent         or space out
        your input 
            text in 
                any way you like to highlight the structure
        of your manuscript and make it easier to edit.
LaTeX fills lines and adjusts spacing between words to produce an
aesthetically pleasing result.

Completely blank lines in the input file break your text into
paragraphs.
To change the font for a single character, word, or set of words, 
enclose the word and the font changing command within braces, 
{\em like this}.
A font changing command not enclosed in braces, like the change to \bf 
bold here, keeps that change in effect until the end of the document or
until countermanded by another font switch, like this change back to 
\rm roman.

\subsection {Displayed Text}
Use the ``quote'' and  ``quotation'' environments for typesetting quoted
material or any other text that should be slightly indented and set off
from the normal text.
\begin{quotation}
The quote and quotation environments are similar, but use different 
settings for paragraph indentation and spacing.

\em When in doubt, consult the manual.
\end{quotation}


\begin{enumerate}
\item 
The ``enumerate'' environment numbers the list elements, like this.

Items in a list can contain multiple paragraphs.
These paragraphs are appropriately spaced and indented according to their
position in the list.
   \begin{itemize}
   \item The ``itemize'' environment sets off list items with ``bullets'',
like this.  Finally, the ``description'' environment lets you put your own
      \begin{description}
      \item[A] label on each item, like this ``A''.
      \item[If the label is long,] the first line of the item text will
be spaced over to the right as needed.
      \end{description}
   \item Of course, lists can be nested, each type up to at least four levels.
One type of list can be nested within another type.
      \begin{itemize}
      \item Nested lists of the same type will change style of numbering 
or ``bullets'' as needed.
      \end{itemize}
   \end{itemize}
\item Don't forget to close off all list environments with the 
appropriate \verb+\end{...}+ command.
Indenting \verb+\begin{...}+, \verb+\item+, and \verb+\end{...}+
commands in the input document according to their nesting level can help 
clarify the structure.
\end{enumerate}

Here is a very simple table showing data lined up in columns.
Notice that the table is in a ``center'' environment to display
it properly.
The title is created simply as another paragraph in the center environment,
rather than as part of the table itself.
\begin{center}
Numbers of Computers Network, By Type.

\begin{tabular}{lr}
Macintosh&175\\
DOS/Windows PC&60\\
UNIX Workstation or server&110\\
\end{tabular}
\end{center}

Here is a more complicated table that has been boxed up, with a multi-column
header and paragraph entries set in one of the columns.
\begin{center}
\begin{tabular}{|l|c|p{3.5in}|}
\hline
\multicolumn{3}{|c|}{Places to Go Backpacking}\\ \hline
Name&Driving Time&Notes\\
&(hours)&\\ \hline
Big Basin&1.5&Very nice overnight to Berry Creek Falls from
either Headquarters or ocean side.\\ \hline
Sunol&1&Technicolor green in the spring.  Watch out for the cows.\\ \hline
Henry Coe&1.5&Large wilderness nearby suitable for multi-day treks.\\ \hline
\end{tabular}

\subsection {Mathematical Equations}
Simple equations, like $x^y$ or $x_n = \sqrt{a + b}$ can be typeset right
in the text line by enclosing them in a pair of single dollar sign symbols.

A more complicated equation should be typeset in {\em displayed math\/} mode,
like this:
\[
p(x)=\lim_{N \rightarrow \infty, \Delta x \rightarrow 0}\sum_{j=1}^{K}(f_j/\Delta x)
\]
The ``equation'' environment displays your equations, and automatically
numbers them consecutively within your document, like this:
\begin{equation}
\mu_m=\sum_{i=0}^m(-1)^i{m \choose i}\mu_1^i\mu_{m-i}^{\prime} \; \; {\rm where} \; \; {m \choose i}={m! \over i!(m-i)!}.
\end{equation}

\end{verbatim}
\section{Including figures}
\LaTeX \ can be used to include figures in documents provided that
they are in {\em Encapsulated} PostScript (EPS) format (ordinary PostScript
is not suitable). The process of including figures can be quite
fiddly and if you need to include a large number of figures it
may be quicker and easier to use a desktop publishing package.
Generally line drawings work better than pictures (i.e. raster
or bitmapped graphics) as the latter take up considerable disk space
and take long periods of time to print. There are different methods
which can be used to include figures. One is to place the 
\verb+\usepackage{graphics}+ command at the beginning of the file
and then use the 
\verb+\includegraphics+ command to include
the EPS file as has been employed here:

\begin{figure}[htp]
\centering
\includegraphics{ulgrid.eps}
\caption{The University of Liverpool Grid (ULGrid) Logo}
\label{fig:logo}
\end{figure}

\section{Keep reference information upto date }

In a document as large and as complicated as a thesis the ability to
cross-reference information contained both within the thesis itself
and in external references quickly and easily is obviously essential. 
Clearly the author does not want to have to change the reference 
numbering every time a reference is added or removed. Fortunately
\LaTeX \ can ensure that the numbering is automatically kept
consistent. References can be cross-references to other information
in the thesis in the form of equations, figures, theorems etc
or they can be references to entries in the bibliography. \LaTeX \ can
also ensure that the tables of contents and figures reflect the current
page numbering. It can also produce a glossary and index automatically. All
of the topics are described in more detail below.

\subsection{Cross-references}
To create
a cross-reference, the 
\verb+\label+ command is used to provide the information
with a key which \LaTeX \ can then use to generate a cross-reference 
number for using the
\verb+\ref+ command. For example, this should
contain the equation number for the second equation displayed above (\ref{eq:summation1}).
The \verb+\label+ command was used earlier to provide a key for the equation
and the \verb+\ref+ command used here to generate the corresponding equation number.
You can extend this idea to figures (e.g. the ULGrid logo: Figure \ref{fig:logo})
or hypotheses, theorems, lemmas, corollaries \ldots infact just about
anything you may want to cross-reference. Note that when you first run
{\bf latex/dvips} on the document in which a reference has been created, the
reference number may not appear correctly. This is because \LaTeX \ cannot
possibly know the number of a reference occuring later in the document to
where it is was cited. To do this, it must scan through the original TeX
file twice and, indeed, the way to ensure that cross-references are upto date
and consistent is just to run {\bf latex} twice before running {\bf dvips}
\footnote{Those from a computing background may see that this is analagous
to how a two--pass compiler resolves forward references.}.

\subsection{Creating a table of contents}
A table of contents can be generated in \LaTeX \ simply by including
the command \verb+\tableofcontents+. As with all references, it may be
necessary to run {\bf latex} twice in order to generate the numbering
correctly. This has been used in this thesis template.

\subsection{Creating a table of figures}
A table of figures can similarly be generated in \LaTeX \ by simply including
the command \verb+\listoffigures+. As with all references, it may be
necessary to run {\bf latex} twice in order to produce the correct numbering.
This has been used in this thesis template.

\subsection{Creating a glossary}
A glossary can be created in a similar manner to an index. The 
\verb+\makeglossary+ is placed in the preamble to document and
\verb+\nomenclature+ used instead of \verb+\index+ to add
glossary entries. The \verb+\printglossary+ command is to indicate
where the glossary is to appear. Creating the index files for a glossary
is slightly more complicated. For this template file the following command
is used:

\begin{verbatim}
$ makeindex -s nomencl -o thesis.gls thesis.glo
\end{verbatim}

The term ``glossary"\nomenclature{{\bf Glossary}}{A list of often difficult or specialized words with their
definitions usually placed in technical literature.} has been included as an example.

\subsection{Creating an index}
Although indices \index{index} \index{indices|see{index}} are difficuly to produce and are
fairly rare in theses, \LaTeX \ can take
out much of
the tedious work needed in creating one. To create
an index\index{index!\LaTeX \ commands} you will need to include the commands \verb+\usepackage{makeindex}+
and \verb+\makeindex+ at the start of the document. The command \verb+\printindex+
is used to indicate where you want the index to appear (usually right at the
end of the document). To index a particular term place the command \verb+\index+
as close as possible to it in the text. You can find a number of references to
this section in the index at the end of this template. Creating the index
requires a little additional work when running \LaTeX. After running the
{\bf latex} command to compile the document, then run\index{index!using makeindex}

\begin{verbatim}
$ makeindex thesis
\end{verbatim}

and run {\bf latex} again twice in order to include the index in the output
({\bf .dvi}) file. Further details can be found in Lamport's book \cite{latex}.


\subsection{Bibliography and citation}
A citation is a cross-reference to another publication, such as a journal
article called the {\em source}. With \LaTeX \ , the citation is produced 
by a \verb+\cite+ command having the citation key as its argument. \LaTeX
uses a separate program called BIBTeX to produce the source list for the
program. You can find details of how to do this in Lamport's book \cite{latex}.
Sources are placed in a separate bibliographic database file ({\bf .bib} file).
The bibliography can be formatted using a variety of predefined styles or you may
be able to find a style sheet that more accurately reflects the format
used in your particular disciple. With BIBTeX it is first necessary to
run {\bf latex} on the main document, then {\bf bibtex} on the {\bf .bib} file
then {\bf latex} twice on the main document again. It may be worth keeping a record
of all references you come across in the {\bf .bib} file for future use as it
can be used in many different publications (e.g. journal papers).

\section{Footnotes}
Footnotes can be added quite easily by using the \verb+\footnote+ command
such as in this example \footnote{There are several footnotes in this template.}.
The University rules state that footnotes should be labelled consecutively
throughout the document. This is what \LaTeX \ does by default.

\section{Splitting the input}

\LaTeX \ uses sophisticated algorithms to decide when to split lines
and how much space to insert between words, essentially mimicking the, now 
bygone, art of manual type setting (composition). Such attention
to detail ensures that the quality of \LaTeX \ output is unsurpassed but it
also comes at a cost. Compiling a \LaTeX \ document, particularly one as large
as a thesis, takes a considerable time which isn't incurred by WYSIWIG
\nomenclature{{\bf WYSIWIG}}{An acronymn for {\em `What you see is what you get'} often
applied to word processors although many in the \TeX \ community believe
that {\em `What you see is all you've got'} is closer to the truth.}
word processors such as Microsoft Word (which consequently produce far inferior 
quality output). Since only small changes are generally made to the document
at a time it makes sense to split the document into smaller parts and compile
these individually using the {\bf latex} command only when they need to be. 
Splitting the input also means
that the author can edit much smaller individual files rather than 
an entire tome\nomenclature{{\bf Tome}}{A book, especially a large or scholarly one.} each time.

The \LaTeX command
\verb+\include+ can be used to gather the sections together in a ``root"
document and this has been used to split Chapter 2 of this template into
a separate file ({\bf chapter2.tex}).  For a thesis it, may
well be worth working on an individual chapter at a time and then combining
them together at the end. The root document will need to be compiled
and any indexes created to ensure the correct numbering before printing
out the final hard copy version. 

\chapter{Thesis Preparation Rules}

The following notes have been produced for the guidance of research 
degree candidates in the presentation of their theses. This information is 
taken from the Postgraduate Handbook and Graduate School web pages. 
To download a PDF version of the Handbook, follow the link on the 
quick links panel of the Graduate School web page at:
\\
\\
http://www.liv.ac.uk/gradschool/thesisprep.htm
\\
\\
or open 
\\
\\
http://www.liv.ac.uk/sas/administration/pgrhandbook.pdf
\\
\\
in Internet 
Explorer.  All students, however, should ensure that they also consult 
their supervisor(s) about the presentation of their theses.

\section{Sources}
Candidates must state generally in the preface and specifically in the body 
of the thesis the sources from which their information is derived and the 
extent to which they have availed themselves of the work of others.

\section{Length}
The various Regulations require that a thesis should be as concise as 
possible. In no circumstances should a thesis of more than 100,000 words 
be submitted for PhD or MD (60,000 for MPhil), including the footnotes 
and appendices, unless written permission has been obtained from the 
candidate's supervisor, the Head of the Department and the Faculty Board 
concerned. It is recommended that the Head of Department seek the view 
of potential examiners before granting his/her approval.

\section{References}
References to published work should be given consistently in a format that 
is currently accepted in the field of work covered by the thesis. If in doubt, 
candidates should consult their supervisors about the best method.

\section{Number of copies required}
Three copies of the thesis and any supporting papers are normally required. 
Two copies should be deposited in the office of the Dean of the candidate's 
Faculty, who will make arrangements for the examination of the thesis. 
The candidate should retain one copy. However, if three Examiners are 
appointed, three copies must be deposited in the office of the Dean of the 
candidate's Faculty.
After the examination, the University will retain two copies in permanent 
bindings. One of these copies will be retained in the University Library. 
The candidate's supervisor will retain the other copy.

\section{Restrictions on access to theses}
An author may impose restrictions on access to theses and copying 
annually for up to five years, if the Head of Department endorses the 
author's statement that such restriction is necessary for good reasons, e.g. 
preparation for publication or a patent application. This will not prevent the 
publication of the Abstract. Permanent restriction is not permitted, nor does 
the University accept theses written under contracts of secrecy (see section 
18 above and the Note on Theses in the current edition of the University 
Calendar). 

\section{Presentation and layout}
In the following specification some of the requirements of BS 4821:1990 
have been adopted to ensure that doctoral theses conform to the standards 
expected by the British Library. Copies of the British Standard (now 
withdrawn from publication) may be consulted in the Sydney Jones and 
Harold Cohen Libraries. Authors' rights are protected under the 
University's agreement with the British Library.
(vii) Typing, printing and copying
Type must be uniform and clear in all copies, for both text and illustrations. 
The minimum height for capital letters is 2 mm and the minimum x-height 
(height of lower-case 'x') 1.5 mm. The main body of the text must be in 
black ink on white paper.
A personal computer with a printer of good quality (e.g. laser or inkjet) 
must be used to produce the first copy. Good, permanent photocopies on 
plain paper are acceptable for the second and third copies. Copies made by 
chemical means, which may fade, are not. The copier must be checked 
before use to ensure that it does not produce extraneous marks on the 
copies.

\section{Binding and lettering}
Theses may be presented for examination in either permanent or temporary 
bindings. 
\begin{description}
\item[Permanent binding]
The thesis to be bound in book form in a strong cloth of any suitable 
colour. Maximum thickness 65 mm (2.5"): if of greater thickness, two or 
more volumes per copy will be required. The binding of all volumes must 
be identical. The thesis should be bound in such a way that it can be 
opened fully for ease of microfilming. Final hardbinding is undertaken off 
university campus by SRJ Ltd in Liverpool (phone: 0151 709 1354). Visit 
their website: http://www.srjservices.co.uk/.
Lettering on permanent bindings to be in gold. Front cover: title of thesis. 
Spine: Top: degree. Middle: surname and initials. Bottom: year of 
submission.
\item[Temporary binding]
The thesis should be presented in such a way that the pages cannot be 
readily removed. The use of ring binders is therefore not permitted. The 
candidate's surname, initials, the date (month and year) and the degree to 
be shown on the outside front cover. Softbinding of initial submission can 
be undertaken by the university print unit.
After the thesis has been approved by the Examiners, two copies must be 
permanently bound as above and deposited with the Dean of the 
candidate's Faculty before arrangements for the conferment of the degree 
can be made.
\end{description}

\section{Title page}
(Centred) Title of thesis
Then `Thesis submitted in accordance with the requirements of the 
University of Liverpool for the degree of Doctor in Philosophy (or other 
degree as appropriate) by $<$full forenames and surname$>$.'
Then (centred) Date (month and year) with suitable line spacing.

\section{Table of contents}
The table of contents must show chapter headings and page numbers. All 
separate sections of the thesis, such as bibliography, lists of abbreviations, 
supporting papers, etc., must also be identified on the contents page.

\section{Abstract}
Each copy of the thesis must be accompanied by a separate copy of the 
Abstract indicating the aims of the investigation and the results achieved. 
For microfilming purposes it must:

\begin{itemize}
\item	Be typed or printed although good photocopies are acceptable;
\item	be not longer than can be accomplished by single-spaced type on 
one side of an A4 sheet (about 450 words);
\item	show the author and title of the thesis in the form of a heading.
\end{itemize}

\section{Paper}
A4 white bond paper of 70 to 100 $g/m^2$ weight must be used for both 
originals and photocopies, except for any endpapers which carry no text. If 
both sides of the paper are used for text, then:

\begin{itemize}
\item	Both sides must be used in both copies which are to be 
permanently bound;
\item	there must be little or no `show-through' - paper lighter than 80 
$g/m^2$ should not be used;
\item	the full binding margin of 40 mm must be allowed on the left 
side of odd pages and the right side of even pages - other 
margins must be 25mm minimum.
\end{itemize}

 Margins and line spacing
1$1 \over 2$ spacing is advised, but at least double line spacing should be used for 
text that contains many subscripts and superscripts. Quotations may be 
indented. Authors should check the text carefully for `widows and orphans' 
and make full use of all error-checking facilities.
\section{Page numbers}
Pages should be numbered consecutively and the position of page numbers 
(candidate's choice or as advised by the supervisor) should be consistent 
throughout.
\section{Footnotes}
Footnotes should be inserted at the foot of the relevant page in single line 
spacing. Smaller type may be used, if available. A line should be ruled 
between footnotes and the text. Footnotes should be numbered 
consecutively throughout the thesis.
\section{Diagrams, maps, illustrations and supporting material}
Diagrams, maps and illustrations should be placed as near to the relevant 
text as possible. If it is necessary to place illustrations in a separate volume, 
the binding must match that of the text. Photographs must be prints of good 
quality and adequate size. Identical and permanent prints of any 
monochrome or colour photographs used must be securely mounted in 
each copy of the thesis.
Published papers submitted in support of the thesis should be sewn in by 
the bookbinder as an appendix.
Essential material that cannot be sewn in (large charts, tapes, floppy discs, 
CDs, microfiches, etc.) must be placed securely in a pocket attached to the 
inside back cover of each copy by the bookbinder. Before submitting 
material that cannot be read without special facilities, candidates must 
satisfy themselves and their supervisors (a) that it is essential to include 
such material and (b) that the Examiners have ready access to such 
facilities. 

\section{Further advice}
The following publications, which can be consulted in the University 
Libraries, give advice on the preparation of theses and methods of 
bibliographical reference. Students are advised to purchase their own 
copies of their chosen manual. 

\subsection{For Humanities and Social Sciences}
\emph{MHRA Style Book}, Modern Humanities Research Association, London. 
MLA Style Sheet, Modern Language Association of America, 
Baltimore. 
\\
\\
Watson, G: \emph{Writing a thesis: A guide to long essays and dissertations}. 
Longman, 1987. 
Turabian, K L: A manual for writers of term papers, theses and 
dissertations. University of Chicago Press, Chicago, 1987. 

\subsection{For Sciences, Engineering and Medicine}
Barrass, R: \emph{Scientists must write: A guide to better writing for 
scientists, engineers and students}. Science Paperbacks, Chapman \& 
Hall, 1978. 
\\
\\
Booth, V: Communicating in science: \emph{Writing and speaking}. 
Cambridge University Press. 
\\
\\
Lindsay, D: \emph{Guide to scientific writing --- a manual for students and 
research workers}. Longman, London, 1990. 
\\
\\
Lock, S: \emph{Thorne's better medical writing}. Pitman, London, 1977. 
O'Connor, M and Woodford, F P: \emph{Writing scientific papers in English: 
An Else-Ciba Foundation guide for authors}. Elsevier, Amsterdam, 
1976. 
\subsection{For all candidates}
Stanisstreet, M: \emph{Writing your thesis: Suggestions for planning and 
writing theses and dissertations in science-based disciplines}. 
University of Liverpool Research Sub-Committee, 1988. 
\\
\\
Stanisstreet, M: \emph{Preparing for your viva: Suggestions for preparing for 
postgraduate vivas in science-based disciplines}. University of 
Liverpool Research Sub-Committee, 1988. 
\\
\\
These booklets were written by a scientist with scientists principally in 
mind, but much of the advice therein will benefit those in other disciplines. 
Copies are normally issued automatically to research students in science-
based departments. They may also be obtained on request, free of charge, 
from Faculty Offices and the Student \& Examinations Division, Senate 
House. 



\appendix
\chapter{Example Appendix}
Appendices are usually labelled with letters separate to ordinary chapters.


\bibliographystyle{plain}
\bibliography{thesis}
\addcontentsline{toc}{chapter}{Bibliography}

\printindex
\addcontentsline{toc}{chapter}{Index}

\begin{verbatim}










\end{verbatim}
\end{document}
